% ==================================================================
% id:       cond-012
% topic:    הסתברות מותנית ואי-תלות
% subtopic: נוסחת ההסתברות השלמה ונוסחת בייס
% solution: full
% ==================================================================
60\% מהמתגייסים לא מתאימים להיות מפקדים, 30\% מתאימים להיות מש״קים ו-10\% מתאימים לקצונה. בבדיקת התאמה לקצונה עוברים בהצלחה 10\% מהלא-מתאימים, 30\% מהמתאימים למש״קים ו-80\% מהמתאימים לקצונה.
\begin{enumerate}
\item מה ההסתברות שמתגייס שנבחר באקראי יעבור את הבדיקה?
\begin{solbox}
נסמן:
$B_1$ = ״המתגייס לא מתאים להיות מפקד״, $B_2$ = ״מתאים להיות מש״ק״, $B_3$ = ״מתאים לקצונה״, ו-$E$ = ״עבר את הבדיקה״.
כל מתגייס שייך בדיוק לאחת משלוש הקבוצות, ולכן $B_1,B_2,B_3$ הם חלוקה של $\Omega$. הנתונים:
\[
\PP(B_1)=0.6,\quad \PP(B_2)=0.3,\quad \PP(B_3)=0.1,\qquad \PP(E\mid B_1)=0.1,\quad \PP(E\mid B_2)=0.3,\quad \PP(E\mid B_3)=0.8.
\]
לפי נוסחת ההסתברות השלמה
\[
\PP(E)=\PP(E\mid B_1)\PP(B_1)+\PP(E\mid B_2)\PP(B_2)+\PP(E\mid B_3)\PP(B_3)=0.1\cdot0.6+0.3\cdot0.3+0.8\cdot0.1=0.06+0.09+0.08=\ans{0.23}.
\]
\end{solbox}
\item מתגייס עבר את הבדיקה. מה ההסתברות שהוא מתאים לקצונה?
\begin{solbox}
מחפשים את $\PP(B_3\mid E)$, כאשר $B_3$ = ״מתאים לקצונה״ ו-$E$ = ״עבר את הבדיקה״. לפי נוסחת בייס ולפי סעיף (א)
\[
\PP(B_3\mid E)=\frac{\PP(E\mid B_3)\PP(B_3)}{\PP(E)}=\frac{0.8\cdot0.1}{0.23}=\frac8{23}\approx\ans{0.348}.
\]
\end{solbox}
\item מתגייס לא עבר את הבדיקה. מה ההסתברות שהוא מתאים להיות מש״ק?
\begin{solbox}
מחפשים את $\PP(B_2\mid E^C)$, כאשר $B_2$ = ״מתאים להיות מש״ק״ ו-$E^C$ = ״לא עבר את הבדיקה״. לפי סעיף (א), $\PP(E^C)=1-0.23=0.77$, ו-$\PP(E^C\mid B_2)=1-0.3=0.7$. לפי נוסחת בייס
\[
\PP(B_2\mid E^C)=\frac{\PP(E^C\mid B_2)\PP(B_2)}{\PP(E^C)}=\frac{0.7\cdot0.3}{0.77}=\frac{0.21}{0.77}=\ans{\frac3{11}}.
\]
\end{solbox}
\end{enumerate}
